% PreTextuais/Abstract.tex
\begin{resumo}[ABSTRACT]
\begin{otherlanguage*}{english}
Public fleet management has historically operated under reactive emergency routines due to fragmented administrative records across asset inventories, maintenance work orders (WOs), and fuel logs. This master's thesis formulates and benchmarks a prescriptive analytics architecture driven by Explainable Artificial Intelligence (XAI), Graph Theory, and Mixed-Integer Linear Programming (MILP). The empirical foundation comprises 77,867 work orders executed between 2020 and 2026 within the State Government of Santa Catarina, totaling BRL 138.76 million in approved and settled expenditures (following BRL 60.05 million retained through technical audit deductions), cross-referenced with 34,280 vehicles and 32.84 million liters of fuel. Nonlinear ensemble learners were evaluated under rolling-window temporal validation. In Canonical Scenario S0 (2,000 vehicles monitored over 36 monthly periods), the MILP model solved by HiGHS with a 0.0000\% optimality gap achieved estimated simulated savings of BRL 9,302,147.26 (17.51\% over the inertial baseline), with 100\% mandatory maintenance compliance. Multimodal data fusion (H1) achieved statistical significance but yielded a marginal error reduction of 0.41\%, failing to reach the pre-registered 15.0\% practical relevance threshold. Explainability benchmarking (H4) revealed a Top-5 Jaccard agreement of 39.29\% between SHAP and LIME, refuting direct method interchangeability and establishing the need for multi-method explanation triangulation in public auditing. Finally, 313 uneconomical vehicles were identified to support statutory disposal procedures. The prescriptive architecture demonstrates measurable efficiency gains and transparent decision-making for public sector governance.

\vspace{\onelineskip}
\noindent
\textbf{Keywords:} prescriptive analytics; public fleet management; graph theory; explainable artificial intelligence; mixed-integer linear programming.
\end{otherlanguage*}
\end{resumo}
